\section{Limitations}

Civilization V offers long-horizon, strategic, multi-agent decision pressure, unique properties close to real-world decision-making scenarios. Yet, it does not reproduce other real-world factors such as command-and-control protocols or institutional constraints. Our outcome measure is the \texttt{use-nuke} flavor, which expresses the strategist's authorization stance, not as a direct launch order. We validated that this proxy reflects genuine authorization intent (Appendix~\ref{app:use-nuke-change-rationales}). The flavor is nonetheless one step removed from nuclear weapon launch in the game.

Because scenarios are drawn from nuke-capable trajectories and high-tension decision points, our analysis focuses on moments where escalation pressure is already present, not whether models would escalate in general (which should surface from emergent self-plays). The replay design captures a single decision point per episode and does not show how models would adapt over subsequent turns.

A downside of using Civilization as a probe, we were unable to reliably convince models that "this is not a game". With all the high-stakes reframing of the prompt, models can still raise "Game Scenario" as a defense, albeit at a lowered ratio (OR 0.51* in the coded corpus). Therefore, the conditions should be only seen as partially effective. That said, had we successfully convinced the models, there is no guarantee that a real-world setting would have a de-escalating effect (see Section \ref{shaping-ethical-reasoning}). 

After our initial tests with a generic ethical prompt produced little impact (Appendix~\ref{app:experimental-conditions}), we added the nuke-specific language, which inevitably confounds the intervention design. That said, the need for this language itself is a finding: most models are so inactive that even with our explicit instruction, only 26.9\% (`ethical x high-stakes') to 48.6\% (`ethical x no-rationale') of reasoning trails contain ethical reasoning keywords (Section~\ref{sec:findings-deductive-reasoning}).

We analyze pre-decision reasoning tokens, yet such tokens may not fully reveal the hidden states that shape the final action \cite{turpin2023unfaithful, lanham2023faithfulness}. We performed the human validation of the keyword tags and the deductive codebook ourselves, and ensemble-versus-human Krippendorff's $\alpha$ reached 0.768 for the codebook, slightly below the conventional 0.8 threshold.

Since most SOTA model providers stopped providing full reasoning tokens, the study only include one (Gemini-3.5-Flash, the strongest Google model as of May 2026) as a sanity check. Our findings from the model was directionally aligned with special caveats (i.e., high-stakes framing can significantly dampen the uptake of ethical prompts), warranting further studies. We call for major providers to open this access to third-party researchers.

None of the interventions or their factorial combinations reliably prevents emergent escalation, and we do not recommend those prompt-based interventions as deployment safeguards for agentic AI in complex, high-stakes settings. Substantial training effort should instead go into models that can reliably and spontaneously surface, engage with, and follow ethical reasoning in decision-making, to the extent that this study's findings become obsolete. The intervention prompts described in Section~\ref{sec:prompt-interventions} and Appendix~\ref{app:experimental-conditions} are documented to support safety research, and we discourage their reuse as production safeguards.
